\section{RESIDUE H.3, H.5, H.6: independently established categorical
content}\label{residue-h.3-h.5-h.6-independently-established-categorical-content}

Research module \texttt{consequences\_criteria}; GPT-6 Astra Pro; effort
setting not exposed; 20 November 2026. The available file
\texttt{sources/INTAKE\_RESIDUE\_20261120.md} reports H.3, H.5 and H.6
at statement level and explicitly says their proofs and cited sources
were not opened by that intake. The underlying \texttt{sections/H.tex}
is absent from the supplied source tree. The results below therefore
have their own complete proofs and precise models; they are not
retrospective certifications of unseen proofs.

\subsection{CC-RES01. Direct image for the universal curve is right Kan
extension}\label{cc-res01.-direct-image-for-the-universal-curve-is-right-kan-extension}

Let \(g,n\geq0\) with \(2g-2+n>0\), and work with the complex analytic
topological Deligne--Mumford stacks \[
X=\overline{\mathcal M}_{g,n+1}^{\mathrm{top}},\qquad
Y=\overline{\mathcal M}_{g,n}^{\mathrm{top}},\qquad
\pi:X\longrightarrow Y.
\] Stratify both by stable dual graph. Let \(\mathcal E\) be a compactly
generated presentable \(\infty\)-category, such as spaces, spectra, or
\(D(R)\) for an associative ring \(R\). Write \(\Pi_X\) and \(\Pi_Y\)
for the actual exit-path \(\infty\)-categories of these stratified
stacks. Then:

\begin{enumerate}
\def\labelenumi{\arabic{enumi}.}
\tightlist
\item
  The forgetful map induces an exit functor
  \(F=\Pi(\pi):\Pi_X\to\Pi_Y\).
\item
  The sheaf direct image \(\pi_*\) takes graph-constructible
  \(\mathcal E\)-valued sheaves to graph-constructible sheaves.
\item
  Under exodromy equivalences \(E_X,E_Y\), there is a natural
  equivalence \[
  E_Y\pi_*\simeq\operatorname{Ran}_{F}E_X.
  \tag{CC-RES.1}
  \] In particular, \[
  (E_Y\pi_*\mathcal F)(y)\simeq
  \lim_{(y\downarrow F)} E_X\mathcal F.
  \tag{CC-RES.2}
  \] The indexing category consists of pairs \((x,\alpha:y\to F(x))\).
  It is not replaced by the strict fibre of \(F\) without a further
  cofinality theorem. For \(\mathcal E=D(R)\), \(\pi_*\) denotes the
  derived sheaf direct image.
\end{enumerate}

\textbf{Inspected primary inputs.} Jansen (2023), Corollary 6.6 and
Theorem 6.7, p.~34, identify these exit categories with Charney--Lee
categories and give their exodromy equivalences; Theorem 6.8, p.~35,
gives the derived functor-category formulation. Jansen (2024), Theorem
2.17, pp.~10--11, gives proper base change for étale topological stacks;
Proposition 3.24, p.~21, gives compactly generated coefficients. Haine,
Porta, and Teyssier (2026 version), Theorem 3.2.3, pp.~34--35, states
that every morphism between exodromic stratified \(\infty\)-topoi
induces the requisite exit functor; Observation 2.3.6, p.~20, identifies
pullback with precomposition. Their Proposition 5.2.11, p.~53, supplies
the product formula used below. These are explicit locators in
arXiv:2308.09551v2, arXiv:2308.09550v3, and arXiv:2401.12825v2
respectively.

\textbf{Complete proof.} Forgetting the last marked point induces a map
of dual-graph strata: remove its leg and stabilize every now-unstable
rational vertex. Contraction of edges, corresponding to smoothing nodes,
is compatible with this operation because stable contraction of families
commutes with base change. Thus \(\pi\) is a stratified map. Jansen's
theorem says that both stratified stacks are exodromic. The
functoriality theorem of Haine--Porta--Teyssier gives \(F\), and its
pullback compatibility gives \[
E_X\pi^*\simeq F^*E_Y,
\] first for space-valued sheaves and then, by the compactly generated
coefficient construction, for \(\mathcal E\)-valued sheaves.

We prove the constructibility claim, which is necessary to identify the
actual sheaf direct image with the right adjoint inside constructible
sheaves. Fix a dual-graph stratum \(Y_G\). On a sufficiently small
contractible analytic chart \(B\to Y_G\), label the finitely many
components, nodes, branches and old marked points. The family over \(B\)
has constant nodal type. Its normalization is a disjoint union of proper
smooth families of compact surfaces, with disjoint labelled sections for
the old marked points and the two preimages of every node. One can
trivialize these families preserving all those sections: choose a
horizontal distribution tangent to the section images, using disjoint
neighbourhoods of the sections and a partition of unity; lift the radial
paths in a small coordinate ball \(B\). Properness makes the lifts exist
throughout these paths after shrinking the ball. They give smooth
trivializations depending continuously on the endpoint. Gluing each pair
of node sections then gives a fibre-preserving homeomorphism \[
\pi^{-1}(B)\simeq C\times B,
\tag{CC-RES.3}
\] where \(C\) is a fixed compact nodal curve, with its nodes and old
marked points distinguished.

This is a stratified trivialization for a refinement of the induced
graph stratification: an extra point in a component's smooth unmarked
part adds a leg to that vertex; at an old marked point it inserts a
rational tail; at a node it inserts a rational bridge. Each of these
conditions is preserved by the trivialization. Thus a
graph-constructible sheaf on \(X\) restricts to a constructible sheaf
for this finite point-and-component stratification on \(C\times B\).

Both \(C\) with this stratification and the ball \(B\) are conically
stratified. The exit product formula gives \[
\Pi(C\times B)\simeq\Pi(C)\times\Pi(B)\simeq\Pi(C),
\] since \(B\) is contractible. Exodromy therefore shows that every such
constructible sheaf on \(C\times B\) is pulled back from a constructible
sheaf \(\mathcal G\) on \(C\). These conical, finite-dimensional spaces
use the usual constructible sheaf categories; the equivalent hypersheaf
formulation of the product theorem yields the same objects here.

The map \(C\to *\) is proper. Proper base change for the product square
now identifies the direct image to \(B\) of the pullback of
\(\mathcal G\) with the constant sheaf having value
\(\Gamma(C,\mathcal G)\). Proper base change for \(\pi\) along
\(Y_G\hookrightarrow Y\), followed by the chart \(B\to Y_G\), identifies
this with \((\pi_*\mathcal F)|_B\). The base-change theorem holds also
for compactly generated coefficients: express \(\mathcal E\)-valued
sheaves through finite-limit-preserving functors from the opposite of the category of compact objects to
space-valued sheaves; evaluation detects the required equivalence,
reducing it to the space-valued statement. Hence \(\pi_*\mathcal F\) is
locally constant on each graph stratum. This proves (2).

It remains to identify the functor. For every
\(A\in\operatorname{Fun}(\Pi_Y,\mathcal E)\) and
\(\mathcal F\in\operatorname{Shv}_{\mathrm{cbl}}(X;\mathcal E)\),
exodromy and the sheaf adjunction give naturally \[
\begin{aligned}
\operatorname{Map}(A,E_Y\pi_*\mathcal F)
&\simeq\operatorname{Map}(E_Y^{-1}A,\pi_*\mathcal F)\\
&\simeq\operatorname{Map}(\pi^*E_Y^{-1}A,\mathcal F)\\
&\simeq\operatorname{Map}(F^*A,E_X\mathcal F).
\end{aligned}
\] Thus \(E_Y\pi_*E_X^{-1}\) is right adjoint to precomposition by
\(F\). By uniqueness of right adjoints, it is right Kan extension. Its
usual pointwise limit formula gives (CC-RES.2). This proves all
assertions. \(\square\)

\textbf{Relation to the intake.} This independently establishes a
precise geometric version of H.6 for the canonical exit categories.
Jansen identifies \(\Pi_{g,n}\) with
\(\mathrm{SC}_{g,n}^{\mathrm{op}}\); consequently the statement can be
expressed using presheaves on entrance categories, with right Kan
extension along the opposite \(u^{\mathrm{op}}\) of the entrance functor
\(u\). What remains uninspected is the equality of the unseen archive's
chosen, explicitly presented forgetful functor with this canonical
functor. The assertion is complex topological. No Frobenius action, Tate
twist, or finite-field weight filtration is created by exodromy; the
separate finite-field statement is CC-ST01--CC-ST04.

\subsection{CC-RES02. Complete lifting-category calculation for the
nodal local
model}\label{cc-res02.-complete-lifting-category-calculation-for-the-nodal-local-model}

Let \(\mathcal C\) be the ordinary category with objects \(0,1\), with
\[
\operatorname{End}(0)=\{1_0\},\quad
\operatorname{End}(1)=\mathbf Z,\quad
\operatorname{Hom}(0,1)=\{a\},\quad
\operatorname{Hom}(1,0)=\varnothing,
\] where composition in \(\mathbf Z\) is addition and \(m\circ a=a\).
Define \[
F:\mathcal C\times\mathcal C\longrightarrow\mathcal C
\] by \(F(1,1)=1\) and \(F(x)=0\) on the other three objects. On the
endomorphisms of \((1,1)\), set \(F(r,s)=r+s\); all arrows between
objects mapped to 0 go to \(1_0\), and all arrows from those objects to
\((1,1)\) go to \(a\).

For \(x\in\mathcal C^2\) and \(\alpha:F(x)\to y\), define the strict
lifting category \(\mathcal L(x,\alpha)\) as follows. Its objects are
arrows \(\beta:x\to z\) such that \(F(z)=y\) and \(F(\beta)=\alpha\). A
morphism from \(\beta\) to \(\beta'\) is an arrow \(\delta:z\to z'\)
with \(F(\delta)=1_y\) and \(\delta\beta=\beta'\). Then \[
|N\mathcal L(x,\alpha)|\simeq
\begin{cases}
S^1,&x=(0,0),\ \alpha=a:0\to1,\\
*,&\text{all other possible pairs }(x,\alpha),
\end{cases}
\tag{CC-RES.4}
\] In the exceptional case the lifting category itself is the one-object
groupoid \(B\mathbf Z\). In the remaining groupoid cases it is
equivalent to the terminal category; the cases with an initial object
need only have contractible nerve.

\textbf{Complete proof.} If \(F(x)=0\) and \(\alpha=1_0\), the lifting
category is the undercategory of \(x\) in the full subcategory
\(F^{-1}(0)\). The identity arrow of \(x\) is initial, so its nerve is
contractible.

Suppose \(F(x)=0\) and \(\alpha=a\). The target must be \(z=(1,1)\). Its
endomorphisms over \(1_1=0\in\mathbf Z\) form \[
K=\ker(\mathbf Z^2\xrightarrow{+}\mathbf Z)
 =\{(r,-r):r\in\mathbf Z\}\simeq\mathbf Z.
\] For \(x=(0,0)\) there is exactly one arrow \(\beta:x\to(1,1)\), and
every element of \(K\) fixes it. Hence the lifting category is exactly
the one-object groupoid \(BK\).

For \(x=(1,0)\) the possible arrows \(\beta\) are indexed by
\(\mathbf Z\), coming from the first coordinate. The element \((r,-r)\)
acts by translation by \(r\). This action is free and transitive, so its
action groupoid is equivalent to the terminal category. For \(x=(0,1)\)
the same argument uses translation by \(-r\) in the second coordinate.

Finally, if \(F(x)=1\), then \(x=(1,1)\), \(y=1\), and \(\alpha\) is an
arbitrary integer \(m\). Its lifts are the pairs \((r,s)\) with
\(r+s=m\). The group \(K\) acts freely and transitively on this set by
addition, so the lifting category again is equivalent to the terminal
category. No arrow \(1\to0\) exists, so these cases exhaust every
possible \((x,\alpha)\). This proves (CC-RES.4). \(\square\)

\textbf{Geometric interpretation with its precise limit.} The exit
category of a complex disk stratified by its origin and punctured disk
is \(\mathcal C\): the punctured stratum has fundamental group
\(\mathbf Z\) and the origin is the cone's initial object. The product
stratification of a bidisk has exit category \(\mathcal C^2\). The
multiplication map \((u,v)\mapsto uv\) induces the displayed functor
because the winding numbers add on \((\mathbf C^*)^2\). Knudsen's
universal-curve local model is exactly \(t=uv\) (1983, proof of Theorem
2.7, p.~179); the analytic nodal coordinates are also used explicitly in
Jansen (2023), Appendix B.2, pp.~43--44. At \((u,v)=(0,0)\) the extra
marked point lies at the node of the contracted curve; stabilization
inserts a rational bridge. Thus the single \(B\mathbf Z\) case is
precisely the local bridge-contraction followed by smoothing described
in the intake. This verifies that local mechanism without asserting a
global equivalence of all the archive's lifting categories with these
local ones.

\subsection{CC-RES03. Exact remaining scope of
H.5}\label{cc-res03.-exact-remaining-scope-of-h.5}

The independent results above resolve the actual sheaf/right-Kan
identification and compute the exceptional nodal lifting category in a
fully specified model. They do not by themselves prove that an arbitrary
categorical base change of the archive's entrance functor preserves
\(W_\infty\)-equivalences, nor its claimed global smoothness
classification. Properness of \(\pi\) supplies geometric proper base
change; that statement quantifies over geometric base changes and cannot
be silently replaced by the universal quantifier over functors in a
definition of \(W_\infty\)-properness. Likewise, a noncontractible local
lifting category does not exclude a possible contraction caused by
additional global morphisms unless the local-to-global comparison is
proved. The missing input is the archive's full functor, the definitions
it uses, and its comparison of global lifting categories with the local
models. H.5 remains an uninspected archive assertion; no contradiction
has been established.

\subsection{References}\label{references}

Haine, P. J., Porta, M., \& Teyssier, J.-B. (2026). \emph{Exodromy
beyond conicality} (arXiv:2401.12825v2; revised 24 June 2026).
https://arxiv.org/abs/2401.12825v2 . Inspected: Observation 2.3.1,
p.~19; Observation 2.3.6, p.~20; Theorem 3.2.3 and its proof,
pp.~34--35; Proposition 5.2.11, p.~53.
https://arxiv.org/pdf/2401.12825v2 .

Jansen, M. Ø. (2023). \emph{Moduli stack of stable curves from a
stratified homotopy viewpoint} (arXiv:2308.09551v2).
https://arxiv.org/abs/2308.09551v2 . Inspected: Corollary 6.6 and
Theorem 6.7, p.~34; Theorem 6.8, p.~35; Appendix B, local coordinates
and trivializations, pp.~40--45. https://arxiv.org/pdf/2308.09551v2 .

Jansen, M. Ø. (2024). Stratified homotopy theory of topological
\(\infty\)-stacks: A toolbox. \emph{Journal of Pure and Applied Algebra,
228}(11), Article 107710. https://doi.org/10.1016/j.jpaa.2024.107710 .
Inspected preprint arXiv:2308.09550v3: Theorem 2.17 and proof,
pp.~10--11; Definition 3.12 and following discussion, pp.~16--17; Remark
3.15, p.~17; Proposition 3.24, p.~21. https://arxiv.org/pdf/2308.09550v3
.

Knudsen, F. F. (1983). The projectivity of the moduli space of stable
curves, II: The stacks \(M_{g,n}\). \emph{Mathematica Scandinavica, 52},
161--199. https://doi.org/10.7146/math.scand.a-12001 . Theorem 2.7 and
its displayed formal local map, p.~179, visually inspected:
https://journals.msp.org/mscand/article/download/1622/1621/1653 .
